% Independent derivation. Prefix SHP. Requires amsmath, amssymb, hyperref.
\section{The full support heat algebra and its exact global-function bridge}
\label{sec:supported-heat-product}

\paragraph{Sources and retained objects.}
The underlying $G_L(R)$ is the existing full-lattice programme
construction, retained under its foundation byline The Clankers in
\href{https://github.com/KokunoYumeto/zeta-function-research-reader/blob/a2851f9eb352e7e4376bc629de86fa4ed9c1c434/workbenches/splitzero-tandem/foundations/split-support-geometry-v11/split_support_geometry_arithmetic_curve_v11.tex}{\emph{Split
Support Geometry: Universal Zero Fibres, Arithmetic Curves, Frobenius
and Perfect Quantization}, v11, Definition \texttt{def:lattice-split}}.
Its supported-zero localization has the public proof
\href{https://github.com/KokunoYumeto/zeta-function-research-reader/blob/3c78cfbd9e194d35117b33277c421c6564ce3dbf/workbenches/splitzero-tandem/continuations/20260923-full-prime-support/031/01_REBUILT_MATHEMATICS.tex}{031,
Theorems \texttt{thm:rebuilt-full-support} and \texttt{thm:e-stalk}}.
The present full support calculation uses HSL1--35 of
\nolinkurl{FULL_LATTICE_HEAT_OPERATOR.tex} for the independently
labelled heat operators, their analytic completion and their
exact coefficient map. The analytic transported-product calculation
is HTP3--10 of \nolinkurl{HEAT_TRANSPORTED_PRODUCT.tex}; its
coefficients and convergence bound are stated at their use below.
The HTP source actually read here is lines 1--256, including
the complete proofs HTP1--12, at SHA--256
\begin{center}\small
\nolinkurl{b81c4489625eed9420eb202009be7b6185cc060bd83f942053b3088d80d91474}.
\end{center}
The actual theta specialization uses THS8--15 and THS18--23
of \nolinkurl{THETA_HEAT_FULL_SUPPORT.tex}, whose human-source
kernel is Brad Rodgers and Terence Tao,
\href{https://arxiv.org/abs/1801.05914v5}{\emph{The De Bruijn--Newman
constant is non-negative}, arXiv:1801.05914v5}.
No zero-location theorem is a premise of the following algebra.

Let $L$ be the entire nontrivial bounded distributive lattice.
Retain
\[
 S=G_L(\mathbb C),\qquad
 (a,\lambda)+(b,\mu)=(a+b,\lambda\vee\mu),\quad
 (a,\lambda)(b,\mu)=(ab,\lambda\wedge\mu).
\]
Thus $a\ne0$ forces $\lambda=1_L$,
$\tau=(0,0_L)$ is the additive zero,
$e=(0,1_L)$ is supported zero, and $(1,1_L)$ is the unit.
The original Taylor coefficients are denoted $a_n$:
\begin{equation}\tag{SHP1}
 \mathcal E=\left\{f(z)=\sum_{n\ge0}a_nz^n:
       \|f\|_h:=\sum_{n\ge0}|a_n|\sqrt{n!}\,h^n<\infty
                                  \text{ for every }h>0\right\}.
\end{equation}
This is the complete entire-function space proved in HSL15--18;
the coefficients themselves have not been changed.
Multiplication of entire functions preserves it. In fact
\begin{equation}\tag{SHP2}
 \|fg\|_h\le\|f\|_{\sqrt2h}\|g\|_{\sqrt2h}.
\end{equation}
For coefficients $a_r,b_s$, use the absolutely convergent
positive sum and
$\sqrt{(r+s)!}\le2^{(r+s)/2}\sqrt{r!s!}$, which follows from
$\binom{r+s}r\le2^{r+s}$. Summing proves (SHP2).
Thus $\mathcal E$ is a commutative unital complex algebra with
its original pointwise function product.

\subsection{Parity-tail closure is an exact multiplicative nucleus}

We first prove the support identity in a setting where every
join written in the proof exists. Put
$\widehat L=\operatorname{Idl}(L)$, with principal embedding
$\jmath(\lambda)=\downarrow\lambda$, ideal intersection as
meet, and ideal generation from unions as arbitrary join.
Finite meets distribute over arbitrary joins: if
$c\in I\cap\bigvee_\alpha J_\alpha$, write
$c\le b_1\vee\cdots\vee b_r$ with each $b_k$ in an indicated
$J_{\alpha_k}$, and then
$c=\bigvee_k(c\wedge b_k)$ is in
$\bigvee_\alpha(I\cap J_\alpha)$.
This proves precisely the distributivity used below.

For sequences $\nu,\mu$ in this complete lattice define
\begin{equation}\tag{SHP3}
 \begin{gathered}
 (\nu*\mu)_n=\bigvee_{r+s=n}\nu_r\wedge\mu_s,\qquad
 (P\nu)_i=\bigvee_{\substack{j\ge i\\j-i\text{ even}}}\nu_j,\\
 \nu+\mu=(\nu_i\vee\mu_i)_{i\ge0},\qquad
 \mathbf u=(1,0,0,\ldots).
 \end{gathered}
\end{equation}
Here the symbols $1,0$ mean the top and bottom of the chosen
lattice. Convolution is associative: both triple products at
degree $n$ are the join of $\nu_r\wedge\mu_s\wedge\omega_k$
over $r+s+k=n$. It has unit $\mathbf u$, is commutative,
and distributes over the coordinate joins, by finite
distributivity. The operator $P$ is extensive, idempotent
and preserves arbitrary joins. Nested same-parity tails prove
these assertions directly.

The needed multiplicative identity is
\begin{equation}\tag{SHP4}
 (P\nu)*(P\mu)\le P(\nu*\mu),\qquad
 P((P\nu)*(P\mu))=P(\nu*\mu).
\end{equation}
To prove the first inequality, distribute the row of its left
side. A term at degree $n$ has indices
$r+s=n$, $a\ge r$, $b\ge s$, with $a-r$ and $b-s$ even.
The term is $\nu_a\wedge\mu_b$.
Then $a+b\ge n$ and $a+b-n$ is even, so that term occurs
in $P(\nu*\mu)_n$. This proves the first inequality after
joining all terms. Applying $P$ gives one inequality in the
second formula; the reverse follows from
$\nu\le P\nu$ and $\mu\le P\mu$.
Thus $P$ is a multiplicative nucleus in the explicit sense
of a join-preserving closure satisfying (SHP4).

For already closed sequences $\nu_i\ge\nu_{i+2}$ and
$\mu_i\ge\mu_{i+2}$, its closed product is completely finite:
\begin{equation}\tag{SHP5}
 (\nu\mathbin{\circledast}\mu)_0
       =(\nu_0\wedge\mu_0)\vee(\nu_1\wedge\mu_1),\qquad
 (\nu\mathbin{\circledast}\mu)_n
       =\bigvee_{r=0}^{n}\nu_r\wedge\mu_{n-r}
                         \quad(n\ge1).
\end{equation}
It equals $P(\nu*\mu)$. Indeed, if $n\ge1$, every term of
$(\nu*\mu)_{n+2}$ has indices $r+s=n+2$, and at least one
of $r,s$ is at least $2$. Lowering that index by $2$
weakly increases its label and supplies a term of
$(\nu*\mu)_n$. Therefore $(\nu*\mu)_{n+2}\le(\nu*\mu)_n$
for $n\ge1$. At $n=0$, the terms of degree $2$ with indices
$(2,0)$ and $(0,2)$ are at most $\nu_0\wedge\mu_0$.
The only remaining term is $\nu_1\wedge\mu_1$.
Higher even degrees are already bounded by degree $2$.
This proves (SHP5), including the one additional mixed-parity
contribution at degree zero. The formulas lie in the
original $L$ when the input labels do.

Let $\mathcal F(L)$ be the set of all such parity flags in
$L^{\mathbb N_0}$. Coordinate join, the zero flag and
$\circledast$ make it a commutative unital idempotent
semiring with unit $\mathbf u$.
For associativity, view it in the complete ideal lattice and use
\[
 P\bigl(P(\nu*\mu)*\omega\bigr)
 =P\bigl((\nu*\mu)*\omega\bigr)
 =P\bigl(\nu*(\mu*\omega)\bigr)
 =P\bigl(\nu*P(\mu*\omega)\bigr),
\]
where (SHP4) applies because $P\omega=\omega$ and $P\nu=\nu$.
Commutativity follows in the same way.
The unit follows from $P\mathbf u=\mathbf u$.
Distributivity follows from join preservation of $P$ and
distributivity of convolution. The principal embedding is
injective, so these equations prove the claims in the
original $L$ as well. No infinite join in $L$ has been assumed.

\subsection{The original independent-coefficient algebra and its flag quotient}

Let $\mathcal D_L$ consist of pairs $(f,\nu)$ with
$f=\sum a_nz^n\in\mathcal E$, $\nu_n\in L$,
$a_n\ne0\Rightarrow\nu_n=1_L$, and every inclusive parity-tail
ideal principal:
\begin{equation}\tag{SHP6}
 \bigvee_{\substack{j\ge i\\j-i\text{ even}}}
       \downarrow\nu_j=\downarrow s_i
                    \qquad\text{for each }i.
\end{equation}
This is the exact original-lattice domain of HSL25.
The generator $s_i$ is a finite join of labels from that
tail, because it belongs to the generated ideal; conversely
such a finite dominating join generates the tail ideal.
Consequently $P\nu=(s_i)$ is well defined in $L$ and is a flag.
The labels $\nu_i$ in $\mathcal D_L$ are otherwise independent.

Addition and ordinary coefficient convolution make $\mathcal D_L$
a commutative unital $S$-algebra.
For closure under convolution, apply (SHP4) in $\widehat L$:
the inclusive tail ideals of $\nu*\mu$ have the explicit
principal generators
\begin{equation}\tag{SHP7}
 g_0=(s_0\wedge t_0)\vee(s_1\wedge t_1),\qquad
 g_n=\bigvee_{r=0}^{n}s_r\wedge t_{n-r}\quad(n\ge1),
\end{equation}
where $(s_i)=P\nu$ and $(t_i)=P\mu$.
This follows because the right side of (SHP4) is the
closed product of those principal tails, calculated by
(SHP5). It proves precisely the required principality.
For addition, tail ideals join and are generated by
$s_i\vee t_i$. Multiplication of all coefficient labels
by a scalar label $\lambda$ gives tail generators
$\lambda\wedge s_i$. The amplitudes stay in $\mathcal E$
by (SHP2); if a product coefficient is nonzero, one of its
finitely many contributing amplitude products is nonzero,
forcing the corresponding support meet and hence the
output label to be $1_L$. Addition and scalar multiplication
satisfy the same carrier test. These arguments prove closure,
and the algebra laws follow from the coefficient operations.

Define the full flag algebra and its product by
\begin{equation}\tag{SHP8}
 \begin{gathered}
 \mathcal A_L=\{(f,\nu):f=\sum a_nz^n\in\mathcal E,\
        \nu\in\mathcal F(L),\
                    a_n\ne0\Rightarrow\nu_n=1_L\},\\
 \mathcal P:\mathcal D_L\longrightarrow\mathcal A_L,\quad
              \mathcal P(f,\nu)=(f,P\nu),\\
 (f,\nu)\circ(g,\mu)=(fg,\nu\circledast\mu).
 \end{gathered}
\end{equation}
A flag has principal tail ideal generated by its first
label in that tail, so $\mathcal A_L$ is exactly the image
of $\mathcal P$ and is contained in $\mathcal D_L$.
The formula for $\circ$ is exactly
$\mathcal P((f,\nu)(g,\mu))$.
Its carrier condition follows either from that construction
or from the nonzero-term argument just used.

The map $\mathcal P$ is an onto unital $S$-algebra homomorphism
from ordinary coefficient convolution in $\mathcal D_L$
to $\circ$ in $\mathcal A_L$.
Addition and the unit are preserved directly.
For multiplication its exact identity is
\begin{equation}\tag{SHP9}
 \mathcal P(vw)=\mathcal P(v)\circ\mathcal P(w),\qquad
 v\sim w\Longleftrightarrow
       f_v=f_w\ \text{and }P\nu^v=P\nu^w .
\end{equation}
This is (SHP4) with the amplitude product unchanged.
Thus the displayed relation is a semiring congruence,
and $\mathcal D_L/\!\sim$ is exactly $\mathcal A_L$.
The algebra laws of $\circ$ follow either from the proved
support semiring laws and the amplitude laws, or from this
surjective homomorphism. The unit has amplitude $1$ and
support $\mathbf u$; the zero has amplitude $0$ and the
zero flag. The scalar $(a,\lambda)\in S$ embeds as its
constant coefficient at degree $0$ and $\tau$ at all
higher degrees. Formula (SHP5) proves this is unital and
acts by the original scalar multiplication.

For the original monomial vector $Z$ representing $z$,
the exact product is
\begin{equation}\tag{SHP10}
 Z\circ Z:
 \quad a_2=1,\ a_i=0\ (i\ne2),\qquad
 \nu_0=\nu_2=1_L,\quad\nu_i=0_L\ (i\ne0,2).
\end{equation}
Its constant coefficient is the supported zero $e$,
not the additive zero $\tau$.
More generally $Z^{\circ n}=\mathcal P(z^n)$ has support
$1_L$ exactly at $i\le n$ of the same parity as $n$.
This follows from (SHP9) and the exact parity-tail of
the single monomial label. It exhibits the mixed
even/odd contribution explicitly rather than dropping
it from multiplication.

\subsection{The transported heat product with every label retained}

On $\mathcal A_L$, HSL10 and HSL18 give the actual heat group
\[
 U_t(f,\nu)=(T_tf,\nu),\qquad
 T_t=e^{-t\partial_z^2},\qquad U_t^{-1}=U_{-t}.
\]
The identity here is the restricted supported-zero heat
operator $U_e$; it is the identity of this flag space.
Both the algebra unit and every embedded scalar are fixed.
Define the transported product
\begin{equation}\tag{SHP11}
 v\star_t w=
 U_t\bigl(U_{-t}v\circ U_{-t}w\bigr)
   =\bigl(T_t((T_{-t}f)(T_{-t}g)),\
                         \nu\circledast\mu\bigr).
\end{equation}
This is well defined by (SHP2) and the heat bounds in HSL17.
For example those estimates give the explicit bound
\[
 \|T_t((T_{-t}f)(T_{-t}g))\|_h
 \le\|f\|_{c(h,t)}\|g\|_{c(h,t)},\qquad
 c(h,t)=\sqrt2\bigl(h+\sqrt{2|t|}\bigr)+\sqrt{2|t|}.
\]
In particular the construction is on the complete analytic
coefficient space, not on an assumed commuting degree
truncation.

The amplitude product is exactly the convergent Wick series
\begin{equation}\tag{SHP12}
 f\star_t^{\mathrm{amp}}g
     =\sum_{n=0}^{\infty}
          \frac{(-2t)^n}{n!}f^{(n)}g^{(n)} .
\end{equation}
The complete analytic proof is HTP7--10 of
\nolinkurl{HEAT_TRANSPORTED_PRODUCT.tex}. Its literal bound,
in the same original-coefficient seminorms, is
\begin{equation}\tag{SHP13}
 \sum_{n\ge0}\frac{(2|t|)^n}{n!}
       \|f^{(n)}g^{(n)}\|_h
 \le
 \frac{\|f\|_{\sqrt{2h^2+a^2}}
       \|g\|_{\sqrt{2h^2+b^2}}}
      {1-2|t|/(ab)}
        \qquad(a,b>0,\ ab>2|t|).
\end{equation}
For clarity, the finite-polynomial identity behind this
formula retains the exact sign and factor $2$.
On $f(x)g(y)$ write $A=\partial_x$ and $B=\partial_y$.
Evaluation at $x=y=z$ turns $\partial_z$ into $A+B$.
The commuting exponentials therefore combine as
$e^{-t(A+B)^2}e^{tA^2}e^{tB^2}=e^{-2tAB}$;
on polynomials all relevant series terminate.
The bound (SHP13), density of polynomials and continuity
of the transported product extend that identity to
all of $\mathcal E$ with absolute convergence.
No support label is obtained by discarding a zero
amplitude term in this series: its support in the
actual lifted algebra is exactly the product already
proved in (SHP5) and (SHP11).

The algebra $(\mathcal A_L,\star_t)$ is commutative,
associative, distributive and unital with the same unit
and embedded scalars as $(\mathcal A_L,\circ)$.
These assertions follow by the invertible transport
definition, using the algebra laws already proved
for $\circ$ and the linearity of $U_t$.
The exact time-changing algebra isomorphisms are
\begin{equation}\tag{SHP14}
 U_t:(\mathcal A_L,\circ)\xrightarrow{\;\sim\;}
                (\mathcal A_L,\star_t),\qquad
 U_s:(\mathcal A_L,\star_t)\xrightarrow{\;\sim\;}
                (\mathcal A_L,\star_{t+s}).
\end{equation}
For the second identity, insert (SHP11) on both sides
and use $U_{-(t+s)}U_s=U_{-t}$. These maps preserve
every original coefficient label.

\subsection{The full coefficient map is a faithful algebra morphism}

Let $K$ be the tropical semifield with zero,
$K_L=K\otimes_{\mathbb B}L$, and
$\iota_L(\lambda)=1_K\otimes\lambda$.
The semiring map $\rho_L=\beta_K\otimes\operatorname{id}_L$
with $\beta_K(a)=1$ for $a\ne0_K$ and $\beta_K(0_K)=0$
satisfies $\rho_L\iota_L=\operatorname{id}_L$.
Here a tropical sum is zero exactly when both entries
are zero, and a tropical product is zero exactly when
a factor is zero; this proves that $\beta_K$ is a
unital semiring homomorphism.

For any commutative additively idempotent semiring $A$,
define $\mathcal F(A)$ to be its parity-decreasing sequences
for the natural order, with the product in (SHP5) written
using addition and multiplication in $A$.
It is a commutative unital semiring.
One precise proof for possibly incomplete $A$ is to
embed it into its ordered-ideal completion, whose joins
and products are the ideals generated by unions and
pairwise products. Multiplication in that completion
distributes over arbitrary joins by the generator
description. The index proof of (SHP4) uses only that
distributivity and monotonicity, so applies there.
For already decreasing sequences, the proof of (SHP5)
uses only finite sums and monotonicity and returns
entries in the original $A$. The principal-ideal
embedding is injective, so the proved algebra laws
descend to $\mathcal F(A)$.
The unit is $(1_A,0_A,0_A,\ldots)$, not the constant
sequence of units.

Coordinatewise application of $\iota_L$ and $\rho_L$
gives unital semiring maps
\begin{equation}\tag{SHP15}
 \mathcal F(L)\xrightarrow{\ \mathcal F(\iota_L)\ }
 \mathcal F(K_L)\xrightarrow{\ \mathcal F(\rho_L)\ }
 \mathcal F(L),\qquad
 \mathcal F(\rho_L)\mathcal F(\iota_L)
                          =\operatorname{id}.
\end{equation}
This follows immediately from the finite formula (SHP5);
order preservation ensures that flags remain flags.
It proves the full-lattice faithfulness of the support
map, including all mixed terms in degree zero.

The complete coefficient map is
\begin{equation}\tag{SHP16}
 \boldsymbol\Phi:\mathcal A_L
     \longrightarrow\mathcal E\times\mathcal F(K_L),
 \qquad(f,\nu)\longmapsto
                  \bigl(f,(\iota_L(\nu_i))_{i\ge0}\bigr).
\end{equation}
It is an injective unital semiring homomorphism for
$\circ$ and ordinary amplitude multiplication.
Both coordinates preserve the operations by (SHP5)
and (SHP15), and equality of their values determines
$f$ and each $\nu_i$. Its image is exactly those
$(f,w)$ for which every $w_i$ belongs to $\iota_L(L)$
and $a_i\ne0$ implies $w_i=1_{K_L}$.
This is the coefficientwise original map
$(a_i,\nu_i)\mapsto(a_i,1_K\otimes\nu_i)$ with its
two coordinates organized into the entire function
and the complete support flag.
For $\star_t$, the same map is a unital homomorphism
to $(\mathcal E,\star_t^{\mathrm{amp}})
\times\mathcal F(K_L)$.
Its heat intertwiner is exactly
$\boldsymbol\Phi U_s=(T_s,\operatorname{id})
\boldsymbol\Phi$, with the products transported
according to (SHP14).

\subsection{Synchronization to the global function carrier is a localization}

The independently labelled algebra is not identified
with one globally labelled function. Its exact morphism
to that object is
\begin{equation}\tag{SHP17}
 g:\mathcal A_L\longrightarrow G_L(\mathcal E),\qquad
 g(f,\nu)=(f,\nu_0\vee\nu_1).
\end{equation}
This belongs to the stated carrier: if $f\ne0$, some
$a_n\ne0$ forces $\nu_n=1_L$, and its parity flag
forces $\nu_0=1_L$ or $\nu_1=1_L$.
The map is unital and additive.
For multiplication (SHP5) gives the exact identity
\begin{equation}\tag{SHP18}
 (\nu\circledast\mu)_0\vee(\nu\circledast\mu)_1
   =(\nu_0\vee\nu_1)\wedge(\mu_0\vee\mu_1).
\end{equation}
Indeed the left side is the join of the four meets
$\nu_i\wedge\mu_j$ for $i,j\in\{0,1\}$, which is
the distributed right side. This proves the
multiplicative assertion, not merely equality of
underlying amplitudes.

There is an explicit zero-preserving, additive and
multiplicative section
\begin{equation}\tag{SHP19}
 j:G_L(\mathcal E)\longrightarrow\mathcal A_L,\qquad
 j(f,\lambda)=(f,(\lambda,\lambda,\lambda,\ldots)),\qquad
 gj=\operatorname{id}.
\end{equation}
If $f\ne0$ its label is top, so the carrier condition
holds at every coefficient; if $f=0$ it imposes
no restriction. The labels form a flag.
Addition is immediate, and (SHP5) shows that the
product of two constant flags is their constant meet.
This proves the three preserved operations and the
section identity.

The section's unit is the following particular idempotent:
\begin{equation}\tag{SHP20}
 J=j(1,1_L)
    =(1,(1_L,1_L,1_L,\ldots)),\qquad
 J\circ J=J,\qquad J\circ v=jg(v).
\end{equation}
It is distinct from the unit of $\mathcal A_L$, whose
support is $(1_L,0_L,0_L,\ldots)$.
For the last identity the amplitude is still $f$.
Its degree-zero support is $\nu_0\vee\nu_1$.
For every $n\ge1$, the support is
$\bigvee_{i=0}^{n}\nu_i=\nu_0\vee\nu_1$,
because the flags decrease separately by parity.
This proves the entire coefficient equality.
Although $J$ has the constant polynomial amplitude $1$,
its support has infinitely many nonzero labels.
It belongs to the completed original-lattice flag
domain and is not a finite-support coefficient vector.
In particular $J$ is not a supported-zero element:
its amplitude is the unit $1$.

For a commutative semiring $A$ and idempotent $c$,
$A[c^{-1}]\simeq cA$ with internal unit $c$.
Indeed an invertible idempotent in a receiving semiring
equals its unit, so every map inverting $c$ factors
uniquely through $a\mapsto ca$; equality there is
exactly $ca=cb$. Applying this proved universal
property to (SHP20) yields
\begin{equation}\tag{SHP21}
 \mathcal A_L[J^{-1}]
       \simeq J\mathcal A_L
       \xrightarrow[\ g\ ]{\ \sim\ }G_L(\mathcal E),
       \qquad g^{-1}=j.
\end{equation}
The map $j$ is unital into the corner whose unit is
$J$, not into all of $\mathcal A_L$.
The localization congruence and its complete fibres are
\begin{equation}\tag{SHP22}
 \begin{gathered}
 (f,\nu)\sim_J(h,\mu)
    \Longleftrightarrow f=h,\
                  \nu_0\vee\nu_1=\mu_0\vee\mu_1,\\
 g^{-1}(f,\lambda)=
 \{(f,\nu):\nu_i\ge\nu_{i+2},\
       \nu_0\vee\nu_1=\lambda,\
       a_i\ne0\Rightarrow\nu_i=1_L\}.
 \end{gathered}
\end{equation}
Necessity is the definition of $g$.
Sufficiency and the complete equality congruence follow
from $Jv=jg(v)$ and injectivity of $j$.
Every displayed target has a preimage (SHP19).
This specifies exactly the mixed-support information
that synchronization loses while retaining the entire
analytic amplitude.

There is a different idempotent that performs amplitude
erasure:
\begin{equation}\tag{SHP23}
 e_0=(0,(1_L,0_L,0_L,\ldots)),\qquad
 e_0\circ(f,\nu)=(0,\nu),\qquad e_0^2=e_0.
\end{equation}
It is precisely the embedded constant scalar $e\in S$.
The support equality follows from (SHP5), since the
only nonzero label of $e_0$ is in degree zero.
Thus
\begin{equation}\tag{SHP24}
 \begin{gathered}
 \mathcal A_L[e_0^{-1}]\simeq\mathcal F(L),\qquad
 \mathcal A_L[(e_0\circ J)^{-1}]\simeq L,\\
 \begin{array}{ccc}
 \mathcal A_L&\xrightarrow{\ g\ }&G_L(\mathcal E)\\
 {\scriptstyle(f,\nu)\mapsto\nu}\big\downarrow
 &&\big\downarrow{\scriptstyle(f,\lambda)\mapsto\lambda}\\
 \mathcal F(L)&\xrightarrow{\ \nu\mapsto\nu_0\vee\nu_1\ }&L .
 \end{array}
 \end{gathered}
\end{equation}
The first isomorphism identifies the corner $e_0\mathcal A_L$
with its full flag and its internal unit with $\mathbf u$.
The product $e_0\circ J$ has amplitude zero and every
coefficient label top. Its corner consists exactly of
zero amplitudes with constant flags, which is $L$.
All maps in the square are unital for their own
displayed target identities and are multiplicative
by (SHP18).
The two routes agree at every element. Localizing
successively at the two commuting idempotents gives
the same corner as localizing at their product:
making their product invertible makes both factors
invertible, and the converse is immediate.
This proves the complete commuting localization
description, including the distinction between
amplitude erasure and support synchronization.

\subsection{The joint coefficient square and its exact localization}

For any commutative idempotent semiring $A$ set
$\epsilon_A:\mathcal F(A)\to A$ by
$\epsilon_A(w)=w_0+w_1$.
The proof of (SHP18), using addition and multiplication
in $A$, proves that this is a unital semiring map.
Its constant-flag section has internal unit
$\mathbf J_A=(1_A,1_A,\ldots)$, and multiplication
by $\mathbf J_A$ sends every flag to the constant
flag of its $\epsilon_A$ value.
The argument uses only the parity inequalities,
so it applies to $A=K_L$.

The full original coefficient map consequently has
the exact commuting square
\begin{equation}\tag{SHP25}
 \begin{array}{ccc}
 \mathcal A_L&
    \xrightarrow{\ \boldsymbol\Phi\ }&
       \mathcal E\times\mathcal F(K_L)\\
 {\scriptstyle g}\big\downarrow
 &&\big\downarrow{\scriptstyle
                   \operatorname{id}_{\mathcal E}\times\epsilon_{K_L}}\\
 G_L(\mathcal E)&
    \xrightarrow{\ \Phi_{\mathcal E}\ }&
       \mathcal E\times K_L,
 \end{array}
 \qquad
 \Phi_{\mathcal E}(f,\lambda)=(f,1_K\otimes\lambda).
\end{equation}
Both composites give
$(f,\iota_L(\nu_0\vee\nu_1))$.
The right vertical map is localization at
$\boldsymbol\Phi(J)=(1,\mathbf J_{K_L})$:
its corner consists of all amplitudes together with
a constant $K_L$-flag, with identity
$(1,\mathbf J_{K_L})$.
The left vertical map is localization at $J$ by
(SHP21). Hence (SHP25) is also a pushout of
unital semirings. Directly, a map from its top-right
object agreeing with a map from the lower-left
object must invert $\boldsymbol\Phi(J)$, and
therefore factors uniquely through the corner
just calculated. The converse is immediate.
This proves the universal property and every
connecting map of the square.

For completely general ideal-labelled flags the same
construction retains a precise completed joint map.
Replace $L$ by $\widehat L$ and $K_L$ by
$\mathfrak K_L=\operatorname{Idl}_{\mathrm{ord}}(K_L)$.
The map on labels is the explicit HSL33 morphism
\[
 \iota^\#(I)=\{w\in K_L:
                   \exists a\in I,\ w\le\iota_L(a)\}.
\]
It preserves arbitrary joins, products and the unit,
and its retraction is the lower ideal generated by
$\rho_L(w)$ for $w$ in the input ideal.
The composition is the identity, because
$\rho_L\iota_L=\operatorname{id}_L$ and both maps
are monotone. Applying these maps to flags preserves
(SHP5), so it gives the injective completed version
of (SHP16) and the complete square (SHP25).
For each original label $\lambda$ this map is
$\downarrow\lambda\mapsto\downarrow\iota_L(\lambda)$.
Thus the completed square restricts to the original
one via explicit principal-ideal embeddings, with
no identification of $L$ with a Boolean lattice.

\subsection{Transported synchronization and the actual theta fibre}

Let $\mathcal E_t$ denote the ring whose underlying
vector space is $\mathcal E$ and whose multiplication
is (SHP12). The unchanged set map $g$ is a unital
semiring homomorphism
$(\mathcal A_L,\star_t)\to G_L(\mathcal E_t)$.
For a map into the ordinary-product global carrier,
the exact morphism instead is
\begin{equation}\tag{SHP26}
 \begin{gathered}
 \gamma_t:(\mathcal A_L,\star_t)\longrightarrow G_L(\mathcal E),
       \qquad \gamma_t(f,\nu)
              =(T_{-t}f,\nu_0\vee\nu_1),\\
 j_t(f,\lambda)=(T_tf,(\lambda,\lambda,\ldots)),\qquad
 \gamma_tj_t=\operatorname{id}.
 \end{gathered}
\end{equation}
The amplitude assertion follows from the transport
definition, which makes $T_{-t}:\mathcal E_t\to\mathcal E$
a unital ring isomorphism. The support assertion is
again (SHP18). Nonzero functions remain nonzero
under $T_t$, so all displayed carriers are respected.

The element $J$ is fixed by every heat map because
$T_t1=1$ and labels are unchanged. It satisfies
\begin{equation}\tag{SHP27}
 J\star_t v=jg(v)=j_t\gamma_t(v),\qquad
 (\mathcal A_L,\star_t)[J^{-1}]
      \simeq G_L(\mathcal E_t)
      \xrightarrow[\ T_{-t}\ ]{\ \sim\ }G_L(\mathcal E).
\end{equation}
For the first equality apply the transport definition
to $J\circ U_{-t}v=jg(U_{-t}v)$ and then apply $U_t$.
Its amplitude returns to $f$ and its label is constant
$\nu_0\vee\nu_1$, proving the equality coefficient by
coefficient. The same idempotent-corner argument
therefore proves the localization statement.
In particular changing the product without changing
the amplitude map has not been silently treated
as a map to the ordinary function ring.

Finally consider the actual theta heat family $H_s$.
THS9 proves that every even Taylor coefficient of
$H_s$ is nonzero at each real $s$, while THS8 gives
all odd coefficients equal to zero.
Consequently the complete fibre of (SHP17) over
$(H_s,1_L)$ at real time is
\begin{equation}\tag{SHP28}
 g^{-1}(H_s,1_L)=
 \{(H_s,\nu):
     \nu_{2j}=1_L\text{ for every }j,\quad
     \nu_1\ge\nu_3\ge\nu_5\ge\cdots
        \text{ is any flag in }L\}.
\end{equation}
Nonvanishing forces each even label to be top.
The odd amplitudes impose no further carrier condition,
and the flag requirement is exactly the displayed
descending odd sequence. Conversely every such
sequence is admissible and has global label top.
This proves both inclusions and thus the full fibre.
The cosine-series lift used in THS18 and HSL31
selects the zero odd flag; it is one specified member
of this complete mixed-support fibre.
The globally constant-label section $j(H_s,1_L)$
instead selects the top odd flag.
Their amplitudes agree, their full support objects
are distinct, and their exact connecting localization
is (SHP21), not an unproved identification.

All results above concern the constructed full support
algebras, their heat products and their exact morphisms.
They keep the actual analytic amplitude throughout
the synchronization map and identify exactly where
amplitude information is discarded by the different
supported-zero localization in (SHP24).
